\chapter{方法与专题导航}\label{app:materials}
本附录按数学结构汇集专题与综合例题链接。可先定位常见信号，再返回定义和定理检查使用条件；章末习题解答见附录\ref{app:answers}。
\section{九章学习路径}
\begin{longtable}{p{36mm}p{88mm}}\toprule 章节&复习主线\\\midrule\endhead
函数、极限与连续&定义与量词、无穷小比较、主项与余项、数列与积分极限。\\
一元函数微分学&求导、微分中值定理、Taylor公式、函数分析与曲率。\\
一元函数积分学&结构整理、换元、分部、有理化、双元、反常与含参积分。\\
常微分方程&分类、降阶、线性结构、特解、定义域与定性分析。\\
向量代数与空间几何&投影、叉积、直线平面、距离夹角与曲面。\\
多元函数微分学&极限与可微、链式法则、隐函数、梯度与极值。\\
重积分&区域描述、换序、坐标变换、对称性与几何应用。\\
曲线与曲面积分&参数化、定向、Green、Gauss、Stokes及势函数。\\
无穷级数&判敛、求和、幂级数、一致收敛与Fourier展开。\\
\bottomrule\end{longtable}

\section{专题方法索引}
\par\noindent\textbf{专题\ref{mat:limits_35}}\quad \hyperref[mat:limits_35:k01]{第1法：定义证明}；\hyperref[mat:limits_35:k02]{第2法：直接代入}；\hyperref[mat:limits_35:k03]{第3法：单侧极限与分段函数}；\hyperref[mat:limits_35:k04]{第4法：概念判断与增长阶}；\hyperref[mat:limits_35:k05]{第5法：子列判别}；\hyperref[mat:limits_35:k06]{第6法：Heine序列化}；\hyperref[mat:limits_35:k07]{第7法：因式分解}；\hyperref[mat:limits_35:k08]{第8法：无穷大转无穷小}；\hyperref[mat:limits_35:k09]{第9法：有理化}；\hyperref[mat:limits_35:k10]{第10法：先求和、求积或恒等改写}；\hyperref[mat:limits_35:k11]{第11法：对数恒等式}；\hyperref[mat:limits_35:k12]{第12法：三角恒等变换}；\hyperref[mat:limits_35:k13]{第13法：两个重要极限}；\hyperref[mat:limits_35:k14]{第14法：一般变量替换}；\hyperref[mat:limits_35:k15]{第15法：等价无穷小}；\hyperref[mat:limits_35:k16]{第16法：洛必达与Stolz差分}；\hyperref[mat:limits_35:k17]{第17法：夹逼准则}；\hyperref[mat:limits_35:k18]{第18法：单调有界与压缩映射}；\hyperref[mat:limits_35:k19]{第19法：微分、积分中值定理}；\hyperref[mat:limits_35:k20]{第20法：Taylor主项}；\hyperref[mat:limits_35:k21]{第21法：导数定义}；\hyperref[mat:limits_35:k22]{第22法：定积分与重积分定义}；\hyperref[mat:limits_35:k23]{第23法：变上限积分极限}；\hyperref[mat:limits_35:k24]{第24法：收敛级数必要条件}；\hyperref[mat:limits_35:k25]{第25法：级数求和求极限}；\hyperref[mat:limits_35:k26]{第26法：常见不等式}；\hyperref[mat:limits_35:k27]{第27法：等价无穷大量}；\hyperref[mat:limits_35:k28]{第28法：Cauchy收敛准则}；\hyperref[mat:limits_35:k29]{第29法：比值极限与根值极限}；\hyperref[mat:limits_35:k30]{第30法：上下极限}；\hyperref[mat:limits_35:k31]{第31法：数学常数识别}；\hyperref[mat:limits_35:k32]{第32法：Wallis公式}；\hyperref[mat:limits_35:k33]{第33法：Stirling公式}；\hyperref[mat:limits_35:k34]{第34法：线性代数方法}；\hyperref[mat:limits_35:k35]{第35法：概率方法与Bernstein逼近}。\par
\par\noindent\textbf{专题\ref{mat:integrals_methods}}\quad \hyperref[mat:integrals_methods:i01]{原函数、常数与分支修正}；\hyperref[mat:integrals_methods:i02]{分项、重组与凑微分}；\hyperref[mat:integrals_methods:i03]{第一换元：内层导数}；\hyperref[mat:integrals_methods:i04]{第二换元：主动参数化}；\hyperref[mat:integrals_methods:i05]{分部积分的选择}；\hyperref[mat:integrals_methods:i06]{多次分部与表格法}；\hyperref[mat:integrals_methods:i07]{循环积分与待定原函数}；\hyperref[mat:integrals_methods:i08]{部分分式：一次与重因子}；\hyperref[mat:integrals_methods:i09]{不可约二次因子}；\hyperref[mat:integrals_methods:i10]{三角奇偶幂}；\hyperref[mat:integrals_methods:i11]{倍角、降幂与积化和差}；\hyperref[mat:integrals_methods:i12]{三角万能代换}；\hyperref[mat:integrals_methods:i13]{正割、余割与递推}；\hyperref[mat:integrals_methods:i14]{双曲函数有理式}；\hyperref[mat:integrals_methods:i15]{二次双元：圆型}；\hyperref[mat:integrals_methods:i16]{二次双元：双曲型}；\hyperref[mat:integrals_methods:i17]{双元降次与递推}；\hyperref[mat:integrals_methods:i18]{分式构型与齐次比值}；\hyperref[mat:integrals_methods:i19]{恒等式诱导与Möbius变换}；\hyperref[mat:integrals_methods:i20]{双元化单元与共同根次}；\hyperref[mat:integrals_methods:i21]{多重根号：和差参数}；\hyperref[mat:integrals_methods:i22]{Euler根式替换}；\hyperref[mat:integrals_methods:i23]{指数单元法}；\hyperref[mat:integrals_methods:i24]{有理积分中的倒数对称}；\hyperref[mat:integrals_methods:i25]{隐函数积分与微分方程}；\hyperref[mat:integrals_methods:i26]{配对与组合积分}；\hyperref[mat:integrals_methods:i27]{含参构造消去对数}；\hyperref[mat:integrals_methods:i28]{复数分解与取实部}；\hyperref[mat:integrals_methods:i29]{伪椭圆构型与恰当代换}；\hyperref[mat:integrals_methods:i30]{三元与椭圆积分}；\hyperref[mat:integrals_methods:i31]{微分方程识别原函数}。\par
\par\noindent\textbf{专题\ref{mat:special_integrals}}\quad \hyperref[mat:special_integrals:erf]{误差函数与Fresnel积分}；\hyperref[mat:special_integrals:sici]{指数、正弦与余弦积分}；\hyperref[mat:special_integrals:polylog]{多重对数与指对双元}；\hyperref[mat:special_integrals:lambert]{Lambert W与隐式变量}；\hyperref[mat:special_integrals:gamma]{Gamma与Beta函数}；\hyperref[mat:special_integrals:bessel]{Bessel函数及加权积分}；\hyperref[mat:special_integrals:airy]{Airy方程与积分}；\hyperref[mat:special_integrals:legendre]{Legendre多项式与正交积分}；\hyperref[mat:special_integrals:hypergeom]{超几何与Mathieu函数}；\hyperref[mat:special_integrals:residue]{Cauchy积分与留数法}；\hyperref[mat:special_integrals:laplace]{Laplace与Fourier积分变换}。\par
\par\noindent\textbf{专题\ref{mat:real_foundations}}\quad \hyperref[mat:real_foundations:complete]{确界、区间套与完备性}；\hyperref[mat:real_foundations:mapping]{映射、反函数与函数方程}；\hyperref[mat:real_foundations:approx]{有理逼近与无理测度}；\hyperref[mat:real_foundations:uniform]{一致连续与反例}。\par
\par\noindent\textbf{专题\ref{mat:differentiation}}\quad \hyperref[mat:differentiation:darboux]{导数的Darboux性质}；\hyperref[mat:differentiation:rolle]{反复Rolle与零点重数}；\hyperref[mat:differentiation:interpolate]{Lagrange插值误差}；\hyperref[mat:differentiation:asymptote]{渐近线与完整函数分析}；\hyperref[mat:differentiation:newton]{Newton法与局部收敛}。\par
\par\noindent\textbf{专题\ref{mat:integral_analysis}}\quad \hyperref[mat:integral_analysis:riemann]{上下和与可积准则}；\hyperref[mat:integral_analysis:primitive]{原函数与可积性不是同一条件}；\hyperref[mat:integral_analysis:holder]{Hölder与Minkowski}；\hyperref[mat:integral_analysis:jensen]{Jensen、Chebyshev与单调权}；\hyperref[mat:integral_analysis:gronwall]{Gronwall积分不等式}；\hyperref[mat:integral_analysis:dirichlet]{反常积分的Dirichlet判别}；\hyperref[mat:integral_analysis:frullani]{Frullani积分}；\hyperref[mat:integral_analysis:numerical]{梯形与Simpson误差}；\hyperref[mat:integral_analysis:parameter]{含参积分的连续与求导}；\hyperref[mat:integral_analysis:applications]{功、流体压力与消费者剩余}。\par
\par\noindent\textbf{专题\ref{mat:geometry}}\quad \hyperref[mat:geometry:rotation]{旋转曲面与轴向识别}；\hyperref[mat:geometry:quadratic]{二次型对角化与主轴}；\hyperref[mat:geometry:curve]{曲线投影与消参}。\par
\par\noindent\textbf{专题\ref{mat:multivariable}}\quad \hyperref[mat:multivariable:topology]{邻域、开闭与紧致性}；\hyperref[mat:multivariable:mixed]{混合偏导交换的条件}；\hyperref[mat:multivariable:jacobian]{反函数与Jacobian链式}；\hyperref[mat:multivariable:homogeneous]{Euler齐次函数定理}；\hyperref[mat:multivariable:constrained2]{条件极值的二阶检验}。\par
\par\noindent\textbf{专题\ref{mat:multiple}}\quad \hyperref[mat:multiple:jacproof]{Jacobian与局部面积倍率}；\hyperref[mat:multiple:tonelli]{Fubini与Tonelli的区别}；\hyperref[mat:multiple:nball]{高维球与Gamma}；\hyperref[mat:multiple:singular]{重积分的奇点判别}。\par
\par\noindent\textbf{专题\ref{mat:fields}}\quad \hyperref[mat:fields:potential]{构造势函数的逐变量积分}；\hyperref[mat:fields:laplacian]{Laplacian与径向场}；\hyperref[mat:fields:curlgrad]{向量恒等式与边界形式}。\par
\par\noindent\textbf{专题\ref{mat:series_methods}}\quad \hyperref[mat:series_methods:condense]{Cauchy凝聚与Raabe判别}；\hyperref[mat:series_methods:rearrange]{绝对收敛、重排与Cauchy积}；\hyperref[mat:series_methods:group]{加括号与原级数}；\hyperref[mat:series_methods:eight]{八种求和入口}；\hyperref[mat:series_methods:odegen]{系数递推与生成函数ODE}；\hyperref[mat:series_methods:abel]{Abel端点定理}；\hyperref[mat:series_methods:diruniform]{一致Dirichlet与逐项求导}；\hyperref[mat:series_methods:weierstrass]{Weierstrass多项式逼近}；\hyperref[mat:series_methods:harmonic]{调和级数的多种发散机制}。\par
\par\noindent\textbf{专题\ref{mat:ode_theory}}\quad \hyperref[mat:ode_theory:direction]{线素场与Euler折线}；\hyperref[mat:ode_theory:riccati]{Riccati方程的已知特解降阶}；\hyperref[mat:ode_theory:clairaut]{Clairaut与d’Alembert型}；\hyperref[mat:ode_theory:trajectory]{正交轨线与模型}；\hyperref[mat:ode_theory:picard]{Picard迭代与存在唯一性证明}；\hyperref[mat:ode_theory:extension]{最大解区间与爆破}；\hyperref[mat:ode_theory:dependence]{连续依赖与Gronwall}；\hyperref[mat:ode_theory:variation]{对初值与参数的可微依赖}。\par
\par\noindent\textbf{专题\ref{mat:ode_systems}}\quad \hyperref[mat:ode_systems:wronskian]{Wronskian与Abel公式}；\hyperref[mat:ode_systems:laplaceode]{Laplace解初值问题}；\hyperref[mat:ode_systems:oscillation]{阻尼、共振与强迫振动}；\hyperref[mat:ode_systems:fundamental]{基本解矩阵与常数变易}；\hyperref[mat:ode_systems:jordan]{矩阵指数与Jordan块}；\hyperref[mat:ode_systems:phase]{相平面、节点、鞍点与焦点}；\hyperref[mat:ode_systems:lyapunov]{Lyapunov稳定与渐近稳定}；\hyperref[mat:ode_systems:cycle]{极限环与径向方程}。\par
\par\noindent\textbf{专题\ref{mat:difference_pde}}\quad \hyperref[mat:difference_pde:difference1]{一阶线性差分}；\hyperref[mat:difference_pde:difference2]{二阶常系数差分}；\hyperref[mat:difference_pde:nonlinear]{非线性迭代与局部稳定}；\hyperref[mat:difference_pde:firstintegral]{首次积分与降维}；\hyperref[mat:difference_pde:transport]{线性齐次PDE与特征线}；\hyperref[mat:difference_pde:quasilinear]{拟线性PDE与特征相交}。\par
\par\noindent\textbf{专题\ref{mat:fourier_advanced}}\quad \hyperref[mat:fourier_advanced:kernel]{Dirichlet核与卷积表达}；\hyperref[mat:fourier_advanced:localization]{局部收敛与Dini条件}；\hyperref[mat:fourier_advanced:fejer]{Fejér核与Cesàro求和}；\hyperref[mat:fourier_advanced:maximal]{Hardy--Littlewood极大函数}；\hyperref[mat:fourier_advanced:weak]{覆盖引理与弱型估计}；\hyperref[mat:fourier_advanced:differentiate]{Lebesgue微分定理}；\hyperref[mat:fourier_advanced:interpolation]{分布函数与Marcinkiewicz插值}；\hyperref[mat:fourier_advanced:banach]{Banach连续性原理}；\hyperref[mat:fourier_advanced:hilbert]{共轭函数与Hilbert变换}；\hyperref[mat:fourier_advanced:hilbertLp]{截断算子与极大Hilbert变换}；\hyperref[mat:fourier_advanced:carleson]{Carleson极大算子与频率调制}；\hyperref[mat:fourier_advanced:basicstep]{基本分解：近场、平均与抵消}；\hyperref[mat:fourier_advanced:dyadic]{二进分割、局部范数与坏集}；\hyperref[mat:fourier_advanced:selection]{频率选择、允许配对与例外集}；\hyperref[mat:fourier_advanced:convergence]{三种收敛层次的边界}；\hyperref[mat:fourier_advanced:rearrangement]{分布函数、递减重排与Lorentz空间}；\hyperref[mat:fourier_advanced:endpoint]{接近L1的空间与Antonov路线}；\hyperref[mat:fourier_advanced:antonov]{有限Fourier系数的集合逼近}。\par
\par\noindent\textbf{专题\ref{mat:fourier_transform}}\quad \hyperref[mat:fourier_transform:gaussian]{Gaussian变换}；\hyperref[mat:fourier_transform:convolution]{卷积、平移与微分}；\hyperref[mat:fourier_transform:plancherel]{反演与Plancherel}；\hyperref[mat:fourier_transform:heat]{热方程与Gaussian核}；\hyperref[mat:fourier_transform:sinecosine]{半轴正弦、余弦变换}；\hyperref[mat:fourier_transform:realmax]{实轴Fourier积分的极大截断}。\par
\par\noindent\textbf{专题\ref{mat:series_sum8}}\quad \hyperref[mat:series_sum8:j1]{一、利用定义求和}；\hyperref[mat:series_sum8:j2]{二、利用已知级数和公式}；\hyperref[mat:series_sum8:j3]{三、利用幂级数求和}；\hyperref[mat:series_sum8:j4]{四、利用Fourier级数求和}；\hyperref[mat:series_sum8:j5]{五、利用代数运算求和}；\hyperref[mat:series_sum8:j6]{六、利用微分方程求和}；\hyperref[mat:series_sum8:j7]{七、利用Euler公式求和}；\hyperref[mat:series_sum8:j8]{八、利用复数求和}。\par
\par\noindent\textbf{专题\ref{mat:advanced_integral_algebra}}\quad \hyperref[mat:advanced_integral_algebra:hermite]{Hermite约化与代数恰当项}；\hyperref[mat:advanced_integral_algebra:elliptic23]{第二、第三类椭圆积分}；\hyperref[mat:advanced_integral_algebra:modbessel]{第二类与修正Bessel积分}；\hyperref[mat:advanced_integral_algebra:mathieu]{Mathieu积分与ODE还原}。\par
\par\noindent\textbf{专题\ref{mat:endpoint_spaces}}\quad \hyperref[mat:endpoint_spaces:qspace]{原子上包络空间Q}；\hyperref[mat:endpoint_spaces:qaspace]{Antonov空间QA}；\hyperref[mat:endpoint_spaces:weakseries]{弱L1的对数加权求和}。\par




\section{综合例题索引}
\par\noindent\textbf{专题\ref{deep:limits_source_deep}}\quad \hyperref[deep:ex:limits_source_deep:eps]{基础例题｜量词与误差估计}；\hyperref[deep:ex:limits_source_deep:heine]{竞赛风格｜Heine定理的双向使用}；\hyperref[deep:ex:limits_source_deep:sin]{考研风格｜同一极限三解}；\hyperref[deep:ex:limits_source_deep:log]{考研风格｜参数相消与分类}；\hyperref[deep:ex:limits_source_deep:exponential]{考研风格｜指数型的二阶精度}；\hyperref[deep:ex:limits_source_deep:rationalization]{考研风格｜根式的三种组织}；\hyperref[deep:ex:limits_source_deep:powersum]{竞赛风格｜幂和的三解}；\hyperref[deep:ex:limits_source_deep:weighted]{竞赛风格｜加权平均原理}；\hyperref[deep:ex:limits_source_deep:factorial]{竞赛风格｜阶乘与对数平均}；\hyperref[deep:ex:limits_source_deep:sqrt]{考研风格｜迭代极限及误差}；\hyperref[deep:ex:limits_source_deep:newton]{竞赛风格｜Newton迭代的二次收敛}；\hyperref[deep:ex:limits_source_deep:singular]{竞赛风格｜带可积奇点的和}；\hyperref[deep:ex:limits_source_deep:laplace]{竞赛风格｜端点集中的积分}；\hyperref[deep:ex:limits_source_deep:gamma]{竞赛风格｜Euler常数存在性}；\hyperref[deep:ex:limits_source_deep:wallis]{竞赛风格｜Wallis乘积证明链}；\hyperref[deep:ex:limits_source_deep:matrix]{竞赛风格｜矩阵迭代与主特征值}；\hyperref[deep:ex:limits_source_deep:probability]{竞赛风格｜Bernstein逼近与概率法}；\hyperref[deep:ex:limits_source_deep:dirichlet]{竞赛拓展｜Dirichlet有理逼近}；\hyperref[deep:ex:limits_source_deep:stirlingcase]{竞赛风格｜阶乘比的三条入口}；\hyperref[deep:ex:limits_source_deep:eirrat]{竞赛拓展｜阶乘尾项与e的无理性}。\par
\par\noindent\textbf{专题\ref{deep:integral_source_deep}}\quad \hyperref[deep:ex:integral_source_deep:dual1]{考研风格｜双元与三角换元双解}；\hyperref[deep:ex:integral_source_deep:dual2]{竞赛风格｜分式构型与共轭因子}；\hyperref[deep:ex:integral_source_deep:dual3]{竞赛风格｜有理化参数的一般公式}；\hyperref[deep:ex:integral_source_deep:euler1]{考研风格｜Euler替换全过程}；\hyperref[deep:ex:integral_source_deep:nested]{竞赛风格｜嵌套根式先降一层}；\hyperref[deep:ex:integral_source_deep:rootratio]{竞赛风格｜根式比值参数}；\hyperref[deep:ex:integral_source_deep:quartic]{考研风格｜四次分母的实因式分解}；\hyperref[deep:ex:integral_source_deep:hermite]{竞赛风格｜Hermite约化的系数方程}；\hyperref[deep:ex:integral_source_deep:binomial]{竞赛风格｜高次双元的可积信号}；\hyperref[deep:ex:integral_source_deep:pseudo1]{竞赛风格｜外观椭圆而实际可换元}；\hyperref[deep:ex:integral_source_deep:pseudo2]{竞赛风格｜有理映射诱导的高次根式}；\hyperref[deep:ex:integral_source_deep:exactalg]{竞赛风格｜由导数反构造积分}；\hyperref[deep:ex:integral_source_deep:triglinear]{考研风格｜一般线性三角分母}；\hyperref[deep:ex:integral_source_deep:trigmix]{竞赛风格｜三角混合与分部闭合}；\hyperref[deep:ex:integral_source_deep:implicit]{竞赛风格｜隐函数参数积分}；\hyperref[deep:ex:integral_source_deep:odeint]{竞赛风格｜由一阶方程生成原函数}；\hyperref[deep:ex:integral_source_deep:lcmroot]{考研风格｜共同根次化有理}；\hyperref[deep:ex:integral_source_deep:expunit]{考研风格｜指数单元与部分分式}；\hyperref[deep:ex:integral_source_deep:cubic]{考研风格｜三次有理积分双解}；\hyperref[deep:ex:integral_source_deep:pellcase]{竞赛风格｜高次恒等式核验}。\par
\par\noindent\textbf{专题\ref{deep:integral_special_deep}}\quad \hyperref[deep:ex:integral_special_deep:arctanparam]{竞赛风格｜Feynman参数法双解}；\hyperref[deep:ex:integral_special_deep:frullani]{竞赛风格｜Frullani公式的严格版本}；\hyperref[deep:ex:integral_special_deep:logsin]{竞赛风格｜对数正弦积分双解}；\hyperref[deep:ex:integral_special_deep:dilog]{竞赛风格｜二重对数反射式}；\hyperref[deep:ex:integral_special_deep:logproduct]{竞赛风格｜对数乘积的完整原函数}；\hyperref[deep:ex:integral_special_deep:elliptic3]{竞赛风格｜三类椭圆标准积分}；\hyperref[deep:ex:integral_special_deep:besselpair]{竞赛拓展｜双Bessel积分与正交}；\hyperref[deep:ex:integral_special_deep:modified]{竞赛拓展｜修正Bessel与Airy积分}；\hyperref[deep:ex:integral_special_deep:legendreproof]{竞赛拓展｜Legendre正交的证明}；\hyperref[deep:ex:integral_special_deep:hypergeneral]{竞赛拓展｜二项式积分的超几何式}；\hyperref[deep:ex:integral_special_deep:mathieu]{竞赛拓展｜Mathieu函数的积分恒等式}；\hyperref[deep:ex:integral_special_deep:incompletebeta]{竞赛拓展｜无穷区间幂权的统一代换}；\hyperref[deep:ex:integral_special_deep:logplus]{竞赛风格｜二重对数的另一类原函数}；\hyperref[deep:ex:integral_special_deep:fresnelprimitive]{竞赛拓展｜Fresnel的尺度与乘积}。\par
\par\noindent\textbf{专题\ref{deep:integral_complex_deep}}\quad \hyperref[deep:ex:integral_complex_deep:cauchy]{竞赛拓展｜高阶极点与Cauchy公式}；\hyperref[deep:ex:integral_complex_deep:residue4]{竞赛拓展｜四次有理积分的留数法}；\hyperref[deep:ex:integral_complex_deep:fourierexp]{竞赛拓展｜指数核、余弦变换与卷积}；\hyperref[deep:ex:integral_complex_deep:laplaceode]{考研风格｜Laplace解受迫初值问题}。\par
\par\noindent\textbf{专题\ref{deep:ode_source_deep}}\quad \hyperref[deep:ex:ode_source_deep:logistic]{考研风格｜分离变量与线性化双解}；\hyperref[deep:ex:ode_source_deep:homogeneous]{考研风格｜齐次方程与几何不变量}；\hyperref[deep:ex:ode_source_deep:bernoulli]{考研风格｜Bernoulli转线性}；\hyperref[deep:ex:ode_source_deep:riccati]{竞赛风格｜Riccati已知特解法}；\hyperref[deep:ex:ode_source_deep:factor]{考研风格｜单变量积分因子}；\hyperref[deep:ex:ode_source_deep:orthogonal]{竞赛风格｜正交轨线与等角轨线}；\hyperref[deep:ex:ode_source_deep:clairaut]{竞赛风格｜Clairaut包络的两条推导}；\hyperref[deep:ex:ode_source_deep:lagrange]{竞赛风格｜Lagrange参数解}；\hyperref[deep:ex:ode_source_deep:nonunique]{竞赛风格｜连续而不唯一}；\hyperref[deep:ex:ode_source_deep:blowup]{竞赛风格｜最大存在区间}；\hyperref[deep:ex:ode_source_deep:sensitivity]{竞赛拓展｜初值可微与变分方程}；\hyperref[deep:ex:ode_source_deep:reduction]{考研风格｜已知解降阶与Euler双解}；\hyperref[deep:ex:ode_source_deep:variation]{考研风格｜常数变易与卷积双解}；\hyperref[deep:ex:ode_source_deep:resonance]{考研风格｜共振三解}；\hyperref[deep:ex:ode_source_deep:complexforce]{考研风格｜复指数与实待定系数}；\hyperref[deep:ex:ode_source_deep:power]{竞赛风格｜幂级数解与显式解互证}；\hyperref[deep:ex:ode_source_deep:generating]{竞赛风格｜生成函数型递推}；\hyperref[deep:ex:ode_source_deep:volterra]{考研风格｜Volterra方程转ODE双解}；\hyperref[deep:ex:ode_source_deep:eulerscheme]{竞赛拓展｜Euler折线的误差分析}；\hyperref[deep:ex:ode_source_deep:periodicA]{竞赛风格｜周期解的充要条件}；\hyperref[deep:ex:ode_source_deep:twosolutions]{考研风格｜两个非齐次特解生成通解}。\par
\par\noindent\textbf{专题\ref{deep:ode_qualitative_deep}}\quad \hyperref[deep:ex:ode_qualitative_deep:jordan]{考研风格｜Jordan与逐次消元双解}；\hyperref[deep:ex:ode_qualitative_deep:rotation]{竞赛风格｜复特征值与极坐标}；\hyperref[deep:ex:ode_qualitative_deep:nonhyperbolic]{竞赛拓展｜线性化失效的两例}；\hyperref[deep:ex:ode_qualitative_deep:limitcycle]{竞赛拓展｜可计算的稳定极限环}；\hyperref[deep:ex:ode_qualitative_deep:diffresonance]{竞赛拓展｜离散共振与生成函数}；\hyperref[deep:ex:ode_qualitative_deep:secondrec]{考研风格｜二阶差分与稳定性}；\hyperref[deep:ex:ode_qualitative_deep:transport]{竞赛拓展｜输运方程的两解}；\hyperref[deep:ex:ode_qualitative_deep:quasilinear]{竞赛拓展｜非线性特征与破裂时间}。\par
\par\noindent\textbf{专题\ref{deep:series_source_deep}}\quad \hyperref[deep:ex:series_source_deep:block]{考研风格｜括号不可逆}；\hyperref[deep:ex:series_source_deep:product]{考研风格｜逐项乘积与Cauchy乘积}；\hyperref[deep:ex:series_source_deep:sparse]{竞赛风格｜收敛不推出$na_n\to0$}；\hyperref[deep:ex:series_source_deep:logseries]{考研风格｜对数级数的两种判别}；\hyperref[deep:ex:series_source_deep:oscillation]{竞赛风格｜振荡级数与端点}；\hyperref[deep:ex:series_source_deep:harmonic4]{基础到竞赛｜调和级数四证}；\hyperref[deep:ex:series_source_deep:harmonicblocks]{竞赛风格｜任意底数与Fibonacci分块}；\hyperref[deep:ex:series_source_deep:harmonictail]{竞赛风格｜尾和的精确极限与一般分块}；\hyperref[deep:ex:series_source_deep:prime]{竞赛拓展｜Erdős素数倒数的计数证明}；\hyperref[deep:ex:series_source_deep:primesecond]{竞赛拓展｜Erdős第二种筛除证明}；\hyperref[deep:ex:series_source_deep:dominated]{竞赛拓展｜控制收敛反证与无界区间}；\hyperref[deep:ex:series_source_deep:means]{竞赛风格｜均值不等式的有限和反证}；\hyperref[deep:ex:series_source_deep:nsquare]{考研风格｜带$n^2$的和三解}；\hyperref[deep:ex:series_source_deep:jdenominator]{考研风格｜分母结构与级数求和}；\hyperref[deep:ex:series_source_deep:uniform1]{考研风格｜一致收敛的域依赖}；\hyperref[deep:ex:series_source_deep:derivative]{竞赛风格｜导数极限不可自动交换}；\hyperref[deep:ex:series_source_deep:endpointtrap]{竞赛风格｜Abel和存在不等于普通收敛}。\par
\par\noindent\textbf{专题\ref{deep:fourier_source_deep}}\quad \hyperref[deep:ex:fourier_source_deep:dini]{竞赛拓展｜Dini收敛准则的完整核推导}；\hyperref[deep:ex:fourier_source_deep:jump]{考研风格｜跳跃点、端点与奇偶延拓}；\hyperref[deep:ex:fourier_source_deep:uniqueness]{竞赛拓展｜Fourier系数唯一性}；\hyperref[deep:ex:fourier_source_deep:differentiation]{竞赛拓展｜Lebesgue微分定理的误差集证明}；\hyperref[deep:ex:fourier_source_deep:hilbertindicator]{竞赛拓展｜示性函数的Hilbert变换}；\hyperref[deep:ex:fourier_source_deep:l2hilbert]{竞赛拓展｜Hilbert变换的$L^2$等距}；\hyperref[deep:ex:fourier_source_deep:carlesongap]{竞赛拓展｜Carleson主估计的逻辑位置}；\hyperref[deep:ex:fourier_source_deep:weaklp]{竞赛拓展｜弱$L^p$与强$L^p$的严格差别}；\hyperref[deep:ex:fourier_source_deep:baselparseval]{竞赛风格｜平方积分求四次倒数和}；\hyperref[deep:ex:fourier_source_deep:threelimits]{竞赛拓展｜三种收敛互不替代}。\par
\par\noindent\textbf{专题\ref{deep:differential_source_deep}}\quad \hyperref[deep:ex:differential_source_deep:logdiff]{考研风格｜对数求导与Taylor互证}；\hyperref[deep:ex:differential_source_deep:implicit2]{考研风格｜隐函数二阶导数}；\hyperref[deep:ex:differential_source_deep:parametric]{考研风格｜摆线曲率与参数导数}；\hyperref[deep:ex:differential_source_deep:leibniz]{考研风格｜高阶导数的有限项化}；\hyperref[deep:ex:differential_source_deep:auxiliary]{竞赛风格｜指数辅助函数双思路}；\hyperref[deep:ex:differential_source_deep:interpolation]{竞赛风格｜插值余项与多点Rolle}；\hyperref[deep:ex:differential_source_deep:logbound]{考研风格｜对数不等式双解}；\hyperref[deep:ex:differential_source_deep:asymptote]{考研风格｜渐近线与精细差值}。\par
\par\noindent\textbf{专题\ref{deep:geometry_source_deep}}\quad \hyperref[deep:ex:geometry_source_deep:skew]{考研风格｜异面直线距离双解}；\hyperref[deep:ex:geometry_source_deep:projection]{考研风格｜点到平面与正交投影}；\hyperref[deep:ex:geometry_source_deep:intersection]{考研风格｜交线切向量的两法}。\par
\par\noindent\textbf{专题\ref{deep:multivariable_source_deep}}\quad \hyperref[deep:ex:multivariable_source_deep:directioncounter]{竞赛风格｜所有方向导数存在仍不可微}；\hyperref[deep:ex:multivariable_source_deep:discontinuous]{竞赛风格｜沿直线都趋零仍无极限}；\hyperref[deep:ex:multivariable_source_deep:implicitmatrix]{考研风格｜二元隐函数的一、二阶信息}；\hyperref[deep:ex:multivariable_source_deep:lagrange]{考研风格｜乘数、代入与不等式三解}；\hyperref[deep:ex:multivariable_source_deep:degenerate]{竞赛风格｜Hessian退化的两个结局}。\par
\par\noindent\textbf{专题\ref{deep:multiple_source_deep}}\quad \hyperref[deep:ex:multiple_source_deep:order]{考研风格｜无初等内层积分的换序}；\hyperref[deep:ex:multiple_source_deep:linearjac]{考研风格｜平行四边形的线性换元}；\hyperref[deep:ex:multiple_source_deep:ellipsoid]{考研风格｜椭球体积与二阶矩}；\hyperref[deep:ex:multiple_source_deep:nball]{竞赛拓展｜高维球体积与Gaussian}；\hyperref[deep:ex:multiple_source_deep:centroid]{考研风格｜半圆薄片的质心}；\hyperref[deep:ex:multiple_source_deep:holder]{竞赛风格｜积分Cauchy与等号条件}；\hyperref[deep:ex:multiple_source_deep:simpson]{竞赛拓展｜求积公式的精度检查}；\hyperref[deep:ex:multiple_source_deep:badlimit]{竞赛风格｜积分与极限失效的尖峰}。\par
\par\noindent\textbf{专题\ref{deep:fields_source_deep}}\quad \hyperref[deep:ex:fields_source_deep:greenarea]{考研风格｜圆周积分的直接与Green法}；\hyperref[deep:ex:fields_source_deep:hole]{竞赛风格｜穿孔区域中的闭而非恰当}；\hyperref[deep:ex:fields_source_deep:hemisphere]{考研风格｜半球通量双解}；\hyperref[deep:ex:fields_source_deep:inverse]{竞赛风格｜反平方场的奇点通量}；\hyperref[deep:ex:fields_source_deep:stokes]{考研风格｜空间环路积分三解接口}。\par
\section{进阶方法与疑难点索引}
\par\noindent\textbf{专题\ref{ic:indefinite_classical_complete}}\quad 
\hyperref[ic:indefinite_classical_complete:normalize]{先整理：除法、配方、拆项与逆链式}；\hyperref[ic:indefinite_classical_complete:logsub]{对数、反函数与整体换元}；\hyperref[ic:indefinite_classical_complete:parts]{分部积分：降阶、升阶与反向选择}；\hyperref[ic:indefinite_classical_complete:cyclic]{循环分部与有限维待定原函数}；\hyperref[ic:indefinite_classical_complete:partial]{部分分式：重根、二次因子与不完全分解}；\hyperref[ic:indefinite_classical_complete:trigparity]{三角有理式的奇偶捷径}；\hyperref[ic:indefinite_classical_complete:trigpowers]{三角幂：留一个微分、降幂与递推}；\hyperref[ic:indefinite_classical_complete:hyperbolic]{双曲有理式与幂递推}；\hyperref[ic:indefinite_classical_complete:eulerthree]{Euler三替换的反解公式}；\hyperref[ic:indefinite_classical_complete:powerrec]{二次根式正负幂的递推}。\par
\par\noindent\textbf{专题\ref{ic:indefinite_dual_complete}}\quad 
\hyperref[ic:indefinite_dual_complete:differentials]{圆型、双曲型与齐次比值}；\hyperref[ic:indefinite_dual_complete:mobius]{恒等式诱导与莫比乌斯变换的构造}；\hyperref[ic:indefinite_dual_complete:reciprocal]{对勾三元与回文多项式}；\hyperref[ic:indefinite_dual_complete:single]{双元化单元与根式和差参数}；\hyperref[ic:indefinite_dual_complete:highdual]{高次双元：统一微分与有理参数}；\hyperref[ic:indefinite_dual_complete:ufamily]{统一有理积分族与公因数降阶}；\hyperref[ic:indefinite_dual_complete:implicit]{隐函数：参数化与约束微分}。\par
\par\noindent\textbf{专题\ref{ic:indefinite_algebra_complete}}\quad 
\hyperref[ic:indefinite_algebra_complete:binomial]{二项微分的三条有理化路线}；\hyperref[ic:indefinite_algebra_complete:hermitemod]{Hermite约化：模多项式算法}；\hyperref[ic:indefinite_algebra_complete:residue]{有理积分的留数式系数法}；\hyperref[ic:indefinite_algebra_complete:quotient]{平方分母的商导数降阶}；\hyperref[ic:indefinite_algebra_complete:triglock]{反三角混合锁与有限维闭合}；\hyperref[ic:indefinite_algebra_complete:explock]{指数锁与归一化变量}；\hyperref[ic:indefinite_algebra_complete:equalparts]{等分锁与对数微分配对}；\hyperref[ic:indefinite_algebra_complete:exactalg]{代数原函数待定与伪椭圆判别}；\hyperref[ic:indefinite_algebra_complete:continuedfraction]{四次根式的多项式连分式算法}；\hyperref[ic:indefinite_algebra_complete:odeidentity]{微分方程引领积分：Lagrange恒等式}。\par
\par\noindent\textbf{专题\ref{ic:indefinite_special_complete}}\quad 
\hyperref[ic:indefinite_special_complete:parameter]{含参不定积分须固定积分常数}；\hyperref[ic:indefinite_special_complete:series]{幂级数逐项积分与命名函数}；\hyperref[ic:indefinite_special_complete:complex]{复指数与实原函数的还原}；\hyperref[ic:indefinite_special_complete:ei]{Ei、li、Si、Ci的识别核}；\hyperref[ic:indefinite_special_complete:gauss]{Gaussian、erf、erfi与Fresnel}；\hyperref[ic:indefinite_special_complete:polylog]{多重对数的逐阶降阶}；\hyperref[ic:indefinite_special_complete:lambert]{Lambert W：用隐式关系换回单元}；\hyperref[ic:indefinite_special_complete:incomplete]{不完全Beta与Gamma的统一幂权}；\hyperref[ic:indefinite_special_complete:elliptic]{三类椭圆标准化与实参数选择}；\hyperref[ic:indefinite_special_complete:boundary]{初等表达的边界与最终验算}。\par
\par\noindent\textbf{专题\ref{ic:delta_series}}\quad 
\hyperref[ic:delta_series:logtest]{对数判别与统一裕量}；\hyperref[ic:delta_series:raabemargin]{Raabe非极限版与临界反例}；\hyperref[ic:delta_series:partialsums]{用部分和作分母的抽象正项级数}。\par
\par\noindent\textbf{专题\ref{ic:delta_limits}}\quad 
\hyperref[ic:delta_limits:riemannerror]{Riemann和的一阶误差}。\par
\par\noindent\textbf{专题\ref{ic:delta_ode}}\quad 
\hyperref[ic:delta_ode:invert]{将自变量与因变量交换}；\hyperref[ic:delta_ode:affine]{仿射分式的平移齐次化}。\par
\par\noindent\textbf{专题\ref{ic:delta_multiple}}\quad 
\hyperref[ic:delta_multiple:permutationpair]{轮换后配成径向函数}；\hyperref[ic:delta_multiple:homogeneousboundary]{齐次边界的径向解析}。\par
\par\noindent\textbf{专题\ref{ic:delta_fields}}\quad 
\hyperref[ic:delta_fields:spaceorientation]{空间交线参数与补线核对}；\hyperref[ic:delta_fields:testfunction]{由积分分部构造零边界测试函数}。\par
\par\noindent\textbf{专题\ref{ic:delta_integral}}\quad 
\hyperref[ic:delta_integral:reciprocalpower]{倒数幂分母的通式与双解}。\par
例题给出题目、思路、详细推导、验算和疑难点；章末习题的完整答案见附录\ref{app:answers}。

\section{结构辨析、巧解与数形结合}
\par\noindent\textbf{第1章}\quad \hyperref[cd:limits-sign]{保号、有限首段与零点定位}；\hyperref[cs:section]{数列的单侧控制、方程列与隐式迭代}；\hyperref[cs:drift]{用可和误差补偿单侧漂移}；\hyperref[cs:roots]{方程列中的统一余项与根的尺度}；\hyperref[cs:inverse]{反向递推与正确的几何缩放}。\par
\par\noindent\textbf{第2章}\quad \hyperref[cd:differential]{微分中的反推、分割与几何辨析}。\par
\par\noindent\textbf{第3章}\quad \hyperref[ci:integrals]{变限积分、临界判敛与几何建模}；\hyperref[ci:oscillatory]{振荡被积函数与平均差商}；\hyperref[ci:parameterdomain]{反常积分参数：先排极点，再谈数值}；\hyperref[ci:kernelwindow]{含绝对值核与移动窗口}；\hyperref[ci:inequality]{积分比较：权函数、累计量和导数包络}；\hyperref[ci:geometry]{分层微元与斜轴旋转}。\par
\par\noindent\textbf{第6章}\quad \hyperref[cmv:section]{偏导条件的减弱、曲线信息与变量选择}；\hyperref[cmv:onepartial]{一个偏导连续的可微判据}；\hyperref[cmv:traces]{两条曲线迹确定全微分的秩条件}；\hyperref[cmv:coordinates]{固定变量不同，偏导算子不同}；\hyperref[cmv:operators]{二阶算子的特征坐标与弯曲路径}。\par
\par\noindent\textbf{第7章}\quad \hyperref[cmi:section]{角向积分、联合分区与未知积分的代数化}；\hyperref[cmi:angular]{极坐标中先积角度}；\hyperref[cmi:partition]{绝对值与最大值的共同分界}；\hyperref[cmi:finiterank]{有限个未知积分化为代数方程}；\hyperref[cmi:regions]{按正负选择积分域与球抛面交域}。\par
\par\noindent\textbf{第8章}\quad \hyperref[cf:section]{曲面度量的消因子、切向消项与逆向识别}；\hyperref[cf:projection]{隐式曲面的投影消因子}；\hyperref[cf:tangent]{切向分量对通量没有贡献}；\hyperref[cf:pathfamily]{路径族的积分最值}；\hyperref[cf:inverse]{恰当条件反求未知场与全局势}。\par
\par\noindent\textbf{第9章}\quad \hyperref[cse:section]{系数运算、隐式项与相邻消去}；\hyperref[cse:radius]{收敛半径的相消与稀疏系数}；\hyperref[cse:pair]{不依赖单调性的相邻消去}；\hyperref[cse:implicit]{有序级数与隐式比较量}；\hyperref[cse:rootphase]{根列与振荡相位定义的级数}；\hyperref[cse:weighted]{离散分部积分保留边界项}。\par
